\chapter{Selected Legislation and Guidance References}
\label{chap:Selected Legislation and Guidance References}

This appendix provides a categorised reference list of the principal legislation, regulations, Approved Codes of Practice, HSE guidance, British Standards, and industry schemes cited throughout \textit{Prudentia Aedificatoria}. It is not a comprehensive legal bibliography; it is a practical reference for construction practitioners seeking to identify the source documents that underpin the assessment frameworks and controls described in the book. All legislation is as in force in England and Wales (and where applicable, Scotland) at the time of writing. Northern Ireland has separate but substantially similar legislation in most areas. Practitioners should verify the current version of any instrument before relying on it in practice.

%---------------------------- SECTION ----------------------------
\section{Primary Legislation}

Primary legislation sets the overarching legal framework from which secondary legislation (regulations) derives its authority. The following Acts are the primary legislative instruments most relevant to construction health and safety in Great Britain.

\begin{table}[H]
\centering
\small
\begin{tabularx}{\textwidth}{>{\raggedright\arraybackslash}p{4.0cm}>{\centering\arraybackslash}p{1.0cm}>{\raggedright\arraybackslash}p{5.0cm}>{\raggedright\arraybackslash}X}
\toprule
\textbf{Title} & \textbf{Year} & \textbf{Scope} & \textbf{Chapters where referenced} \\
\midrule
Health and Safety at Work etc.\ Act & 1974 & The principal UK occupational health and safety statute. Imposes general duties on employers, employees, the self-employed, and those in control of premises. Provides the framework for secondary legislation and the HSE's enforcement powers. & Throughout \\
\midrule
Factories Act & 1961 & Historic predecessor to the HSW Act, now largely superseded but relevant for understanding the development of construction health and safety regulation. & Ch.2 \\
\midrule
Corporate Manslaughter and Corporate Homicide Act & 2007 & Creates a criminal offence for organisations whose gross breach of a duty of care causes a person's death. Applies to companies, partnerships, and public bodies. & Ch.29, Ch.30 \\
\midrule
Health and Safety (Offences) Act & 2008 & Extended the range of health and safety offences triable in the Crown Court and increased maximum penalties for individuals. & Ch.29 \\
\midrule
Equality Act & 2010 & Prohibits discrimination in the workplace based on protected characteristics including disability. Imposes a duty to make reasonable adjustments for disabled workers. Relevant to welfare provision and vulnerable worker risk assessment. & Ch.24, Ch.25 \\
\midrule
Modern Slavery Act & 2015 & Requires organisations with annual turnover $\geq$\pounds36\,m to report on steps taken to ensure modern slavery is not occurring in their operations or supply chains. & Ch.28 \\
\midrule
Building Safety Act & 2022 & Introduced a new regulatory regime for higher-risk buildings (HRBs), including the Building Safety Regulator and the Golden Thread of building information. Directly affects CDM obligations on HRB projects. & Ch.27 \\
\bottomrule
\end{tabularx}
\end{table}

%---------------------------- SECTION ----------------------------
\section{Regulations}

Secondary legislation (Statutory Instruments made under primary Acts) provides the specific, detailed requirements for construction health and safety. The following regulations are those most frequently cited in the chapters of this book.

\begin{table}[H]
\centering
\small
\begin{tabularx}{\textwidth}{>{\raggedright\arraybackslash}p{4.8cm}>{\centering\arraybackslash}p{0.9cm}>{\raggedright\arraybackslash}p{4.2cm}>{\raggedright\arraybackslash}X}
\toprule
\textbf{Title} & \textbf{Year} & \textbf{Scope} & \textbf{Chapters} \\
\midrule
Management of Health and Safety at Work Regulations & 1999 & The management framework: risk assessment, competent persons, emergency procedures, health surveillance, new and expectant mothers, young workers. & Throughout \\
\midrule
Construction (Design and Management) Regulations & 2015 & CDM duty holder framework: client, principal designer, principal contractor, contractor, designer, worker. PCI, CPP, HSF, F10. & Ch.27 \\
\midrule
Provision and Use of Work Equipment Regulations (PUWER) & 1998 & All work equipment: suitability, maintenance, inspection, guarding, training. & Ch.9, Ch.13 \\
\midrule
Lifting Operations and Lifting Equipment Regulations (LOLER) & 1998 & Lifting equipment and accessories: strength, stability, SWL marking, thorough examination, operator competence, planning. & Ch.15 \\
\midrule
Control of Substances Hazardous to Health Regulations (COSHH) & 2002 & Hazardous substances: assessment, prevention, control, health surveillance, emergency procedures. & Ch.11 \\
\midrule
Control of Noise at Work Regulations & 2005 & Noise exposure: action and limit values, noise assessment, hearing protection zones, health surveillance. & Ch.10 \\
\midrule
Control of Vibration at Work Regulations & 2005 & Hand-arm and whole-body vibration: action and limit values, HAV assessment, health surveillance, tiered assessment. & Ch.10 \\
\midrule
Manual Handling Operations Regulations & 1992 & Manual handling: avoidance where reasonably practicable, assessment of unavoidable manual handling, reduction of risk. & Ch.5 \\
\midrule
Personal Protective Equipment at Work Regulations & 2022 & PPE provision, maintenance, storage, training, and use. 2022 amendments extended duties to limb (b) workers. & Ch.23 \\
\midrule
Confined Spaces Regulations & 1997 & Confined space entry: risk assessment, avoidance where reasonably practicable, safe systems of work, rescue arrangements. & Ch.16 \\
\midrule
Work at Height Regulations & 2005 & Work at height: avoidance, prevention, mitigation hierarchy; collective before personal measures; equipment inspection. & Ch.7, Ch.8 \\
\midrule
Regulatory Reform (Fire Safety) Order (RRFSO) & 2005 & Fire safety in non-domestic premises: responsible person duties, fire risk assessment, fire safety arrangements, evacuation. & Ch.14 \\
\midrule
Fire Safety (England) Regulations & 2022 & Additional fire safety duties for high-rise residential buildings following the Grenfell Tower Inquiry recommendations. & Ch.14 \\
\midrule
Control of Asbestos Regulations & 2012 & Asbestos: duty to manage, survey requirements, licensed/notifiable non-licensed/non-notifiable work categories, training. & Ch.12 \\
\midrule
Electricity at Work Regulations & 1989 & Electrical systems: safe design, construction, maintenance, and operation; working on or near live conductors. & Ch.13 \\
\midrule
Gas Safety (Installation and Use) Regulations & 1998 & Gas installations: design, installation, maintenance, and use; emergency procedures for gas strikes. & Ch.6, Ch.22 \\
\midrule
Reporting of Injuries, Diseases and Dangerous Occurrences Regulations (RIDDOR) & 2013 & Reporting of deaths, specified injuries, over-7-day injuries, dangerous occurrences, and occupational diseases. & Ch.29 \\
\midrule
Workplace (Health, Safety and Welfare) Regulations & 1992 & Fixed premises: ventilation, temperature, lighting, cleanliness, welfare. Applies to site offices and welfare units. & Ch.24 \\
\midrule
Working Time Regulations & 1998 & Working hours: 48-hour average maximum; rest periods; annual leave; night work; specific provisions for young workers. & Ch.24, Ch.25 \\
\midrule
Pressure Systems Safety Regulations & 2000 & Pressure systems and transportable pressure vessels: design, installation, safe use, and examination. & Ch.9 \\
\midrule
Control of Lead at Work Regulations & 2002 & Lead exposure: assessment, control, air monitoring, biological monitoring, health surveillance. & Ch.11 \\
\midrule
Ionising Radiations Regulations & 2017 & Ionising radiation: controlled and supervised areas, dose limits, monitoring, local rules. Relevant to radiography on site. & Ch.11 \\
\midrule
Health and Safety (Safety Signs and Signals) Regulations & 1996 & Safety signs: prescribed formats for prohibition, warning, mandatory, and emergency signs; maintenance. & Throughout \\
\bottomrule
\end{tabularx}
\end{table}

%---------------------------- SECTION ----------------------------
\section{Approved Codes of Practice}

Approved Codes of Practice (ACOPs) have a special legal status. Failure to follow an ACOP is not automatically a criminal offence, but it is admissible as evidence of failure to comply with the underlying legislation; the defendant must demonstrate that they complied by other equally effective means. ACOPs are issued by the HSE with the approval of the Secretary of State.

\begin{table}[H]
\centering
\small
\begin{tabularx}{\textwidth}{>{\raggedright\arraybackslash}p{1.2cm}>{\raggedright\arraybackslash}p{4.5cm}>{\centering\arraybackslash}p{0.9cm}>{\raggedright\arraybackslash}p{3.5cm}>{\raggedright\arraybackslash}X}
\toprule
\textbf{Ref.} & \textbf{Title} & \textbf{Year} & \textbf{Scope} & \textbf{Chapters} \\
\midrule
L153 & Managing health and safety in construction (CDM 2015 ACOP) & 2015 & Definitive interpretation of CDM 2015, including duty holder obligations, welfare requirements, and documentation standards. & Ch.24, Ch.27 \\
\midrule
L141 & Safe use of lifting equipment (LOLER ACOP) & 1998 & Lifting equipment: planning, management, thorough examination, operator competence. & Ch.15 \\
\midrule
L140 & Safe use of work equipment (PUWER ACOP) & 1998 & Work equipment: suitability, maintenance, inspection, training. & Ch.9 \\
\midrule
L101 & Safe work in confined spaces (Confined Spaces Regulations ACOP) & 2014 & Confined space risk assessment, safe systems, rescue arrangements. & Ch.16 \\
\midrule
L131 & Workplace health, safety and welfare (Workplace Regs ACOP) & 1992 & Workplace welfare standards for fixed premises. & Ch.24 \\
\midrule
L8 & Legionnaires' disease: The control of legionella bacteria in water systems & 2013 & Water system risk assessment and legionella control. Relevant to welfare units and temporary water systems on construction sites. & Ch.24 \\
\midrule
L5 & Control of substances hazardous to health (COSHH ACOP) & 2013 & COSHH assessment, control measures, health surveillance. & Ch.11 \\
\midrule
L23 & Manual handling: Manual Handling Operations Regulations ACOP & 2016 & Manual handling risk assessment and ergonomic control. & Ch.5 \\
\bottomrule
\end{tabularx}
\end{table}

%---------------------------- SECTION ----------------------------
\section{HSE HSG and L-Series Guidance}

HSE guidance documents (HSG series, L series, and others) are not legally binding but represent the HSE's interpretation of how compliance with the relevant legislation should be achieved. Departure from HSE guidance without equivalent alternative measures is likely to be cited as evidence of inadequate risk management in enforcement proceedings. The documents listed below are those most frequently referenced in the chapters of this book.

\begin{table}[H]
\centering
\small
\begin{tabularx}{\textwidth}{>{\raggedright\arraybackslash}p{1.5cm}>{\raggedright\arraybackslash}p{4.5cm}>{\centering\arraybackslash}p{0.9cm}>{\raggedright\arraybackslash}p{3.2cm}>{\raggedright\arraybackslash}X}
\toprule
\textbf{Ref.} & \textbf{Title} & \textbf{Year} & \textbf{Scope} & \textbf{Chapters} \\
\midrule
HSG47 & Avoiding danger from underground services & 2014 & Service identification, CAT survey, safe digging zones, hand-dig, vacuum excavation. & Ch.6, Ch.22 \\
\midrule
HSG185 & Health and safety in excavations & 2006 & Excavation stability, soil classification, trench support, inspection requirements. & Ch.6 \\
\midrule
HSG150 & Health and safety in construction & 2006 & General construction health and safety guidance: risk assessment, management, welfare. & Throughout \\
\midrule
HSG144 & The safe use of vehicles on construction sites & 2009 & Site traffic management, segregation, banksmen, vehicle selection. & Ch.9, Ch.19 \\
\midrule
HSG33 & Health and safety in roof work & 2012 & Roof work: fragile roofs, edge protection, access equipment, safe access systems. & Ch.7 \\
\midrule
HSG136 & Workplace transport safety & 2014 & Workplace transport risk assessment, traffic management, loading/unloading operations. & Ch.19 \\
\midrule
HSG65 & Managing for health and safety & 2013 & The Plan-Do-Check-Act management framework for health and safety. & Ch.30 \\
\midrule
EH40 & Workplace Exposure Limits (WELs) & 2020 & Maximum workplace concentrations for hazardous substances; legal limits under COSHH. & Ch.11 \\
\midrule
EH44 & Dusty work in construction & 2008 & Respirable dust controls: RPE, LEV, water suppression, enclosure. & Ch.21 \\
\midrule
EH57 & The problems of asbestos remaining in situ in premises & 1999 & Asbestos in buildings: management strategy, risk assessment for maintenance work. & Ch.12 \\
\midrule
L143 & Managing and working with asbestos (Control of Asbestos Regulations ACOP/Guidance) & 2013 & Asbestos: licensed, notifiable non-licensed, and non-notifiable work; training; health surveillance. & Ch.12 \\
\midrule
INDG417 & Leading health and safety at work & 2013 & Board and director responsibilities for health and safety governance. & Ch.30 \\
\midrule
HSG48 & Reducing error and influencing behaviour & 1999 & Human factors in health and safety: error reduction, behavioural approaches. & Ch.30 \\
\midrule
HSG245 & Investigating accidents and incidents & 2004 & Accident investigation methodology: root cause analysis, 5-Whys, fishbone. & Ch.29 \\
\midrule
INDG163 & Five steps to risk assessment & 2014 & The five-step risk assessment method. & Ch.3 \\
\midrule
HSG76 & Warehousing and storage: a guide to health and safety & 2007 & Storage and materials handling; relevant to site materials compound and storage areas. & Ch.9 \\
\bottomrule
\end{tabularx}
\end{table}

%---------------------------- SECTION ----------------------------
\section{British Standards}

British Standards (BS and BS EN) are developed by the British Standards Institution (BSI). They provide detailed technical specifications that, when referenced in legislation, ACOPs, or HSE guidance, become part of the benchmark for compliance. The following standards are the most frequently relevant to the risk domains covered in this book.

\begin{table}[H]
\centering
\small
\begin{tabularx}{\textwidth}{>{\raggedright\arraybackslash}p{2.5cm}>{\raggedright\arraybackslash}p{4.0cm}>{\centering\arraybackslash}p{0.9cm}>{\raggedright\arraybackslash}p{3.0cm}>{\raggedright\arraybackslash}X}
\toprule
\textbf{Reference} & \textbf{Title} & \textbf{Year} & \textbf{Scope} & \textbf{Chapters} \\
\midrule
BS 8484 & Provision of lone worker device services & 2022 & Requirements for lone worker devices, monitoring centres, and emergency response. Compliance enables police grade-1 response. & Ch.26 \\
\midrule
BS 7121 & Code of practice for the safe use of cranes & Multiple parts & Planning, operation, and maintenance of cranes in construction. & Ch.15 \\
\midrule
BS EN 13000 & Cranes --- mobile cranes & 2010 & Design, stability, and safety requirements for mobile cranes. & Ch.15 \\
\midrule
BS EN 280 & Mobile elevating work platforms & 2022 & Design, safety requirements, and testing for MEWPs. & Ch.7 \\
\midrule
BS 6187 & Code of practice for full and partial demolition & 2011 & Demolition planning, methods, and safety management. & Ch.17 \\
\midrule
BS EN 1263 & Safety nets & 2014 & Specification and test methods for safety nets used for fall protection. & Ch.7 \\
\midrule
BS 1722 & Fences --- specification for chain link fences & Multiple parts & Standard for site hoarding and perimeter security fencing. & Ch.26 \\
\midrule
BS EN 13374 & Temporary edge protection systems & 2013 & Requirements for temporary edge protection used during construction. & Ch.7 \\
\midrule
BS EN ISO 45001 & Occupational health and safety management systems & 2018 & Requirements and guidance for OH\&S management systems; replaces OHSAS 18001. & Ch.30 \\
\midrule
BS 5975 & Code of practice for temporary works procedures and the permissible stress design of falsework & 2019 & Temporary works design, procedures, and management. & Ch.8, Ch.18 \\
\midrule
BS EN 12811 & Temporary works equipment --- Scaffolds & 2003 & Performance requirements, design, and use of scaffolding. & Ch.8 \\
\midrule
BS EN 795 & Protection against falls from height --- Anchor devices & 2012 & Requirements for anchor devices used as part of PPE fall arrest systems. & Ch.7 \\
\midrule
BS EN 388 & Protective gloves against mechanical risks & 2016 & Performance requirements and testing for mechanical protective gloves. & Ch.23 \\
\midrule
BS EN 374 & Protective gloves against dangerous chemicals and micro-organisms & 2016 & Performance requirements for chemically protective gloves. & Ch.11, Ch.23 \\
\midrule
BS EN 166 & Personal eye protection --- Specifications & 2002 & Specification for eye and face protection equipment. & Ch.23 \\
\bottomrule
\end{tabularx}
\end{table}

%---------------------------- SECTION ----------------------------
\section{Industry Schemes and Guidance}

The following industry schemes and guidance documents supplement the statutory and standards framework, providing practical tools, training frameworks, pre-qualification processes, and sector-specific guidance that are widely used in UK construction health and safety management.

\begin{table}[H]
\centering
\small
\begin{tabularx}{\textwidth}{>{\raggedright\arraybackslash}p{3.0cm}>{\raggedright\arraybackslash}p{4.5cm}>{\raggedright\arraybackslash}p{3.0cm}>{\raggedright\arraybackslash}X}
\toprule
\textbf{Organisation / Scheme} & \textbf{Description} & \textbf{Key Products / Standards} & \textbf{Chapters} \\
\midrule
SSIP (Safety Schemes in Procurement) & Umbrella body for health and safety pre-qualification schemes. Enables mutual recognition between member schemes. & Core criteria set; mutual recognition framework & Ch.28 \\
\midrule
CHAS (Contractors Health and Safety Assessment Scheme) & SSIP member scheme providing health and safety pre-qualification for contractors. & CHAS certificate; CHAS Advanced; CHAS Premier & Ch.28 \\
\midrule
SafeContractor & SSIP member scheme providing health and safety pre-qualification. & SafeContractor certificate; various tiers & Ch.28 \\
\midrule
Constructionline & SSIP member scheme and supplier information management service for the construction industry. & Bronze, Silver, Gold levels & Ch.28 \\
\midrule
CSCS (Construction Skills Certification Scheme) & Card scheme providing evidence of health and safety training for construction workers and managers. & Green, Blue, Gold, Black card levels & Ch.27 \\
\midrule
CPCS (Construction Plant Competence Scheme) & Competence card scheme for construction plant operators, covering over 60 categories of plant. & CPCS Technical Test; Health, Safety and Environment Test; Competence logbook & Ch.9, Ch.15 \\
\midrule
IPAF (International Powered Access Federation) & Trade association for the powered access (MEWP) sector; administers the PAL card operator training scheme. & PAL card; MEWP categories (1a, 1b, 3a, 3b, 3c, 3d) & Ch.7, Ch.15 \\
\midrule
PASMA (Prefabricated Aluminium Scaffolding Manufacturers Association) & Trade body for the mobile tower scaffold sector; administers operator training and competence cards. & PASMA card; mobile tower scaffold course & Ch.8 \\
\midrule
Mates in Mind & UK charity dedicated to improving mental health in construction; provides training, resources, and the Mates in Mind framework for employers. & Mental Health First Aid; Ambassador training; Helpline & Ch.25 \\
\midrule
CLOCS (Construction Logistics and Cyclist Safety) & Industry standard for construction logistics safety, managing the risk to vulnerable road users from delivery vehicles. & CLOCS Standard; CLOCS compliance audit & Ch.19, Ch.28 \\
\midrule
FORS (Fleet Operator Recognition Scheme) & Voluntary accreditation for vehicle operators: safety, efficiency, and environment standards. & FORS Bronze, Silver, Gold accreditation & Ch.19, Ch.28 \\
\midrule
CITB (Construction Industry Training Board) & Industry training board for construction in Great Britain; funds training and apprenticeships; administers the CSCS Health, Safety and Environment Test. & H\&SE Test; Site Safety Plus courses; SMSTS; SSSTS & Ch.27 \\
\midrule
NHBC (National House-Building Council) & Building standards body and warranty provider for new homes; sets construction standards for housebuilding. & NHBC Standards; Buildmark warranty; technical guidance chapters & Throughout \\
\midrule
CIOB (Chartered Institute of Building) & Professional body for construction management; produces research and guidance on construction health and safety including mental health data. & Mental health research; professional development; competence frameworks & Ch.25 \\
\midrule
CIBSE (Chartered Institution of Building Services Engineers) & Professional body for building services engineers; produces technical guidance for MEP systems relevant to fire and electrical safety. & CIBSE Guides; TM series; AM series & Ch.13, Ch.14 \\
\midrule
Sentencing Council --- Health and Safety Offences Definitive Guideline & Sentencing guidance for health and safety, corporate manslaughter, and food safety offences; determines fines based on culpability, risk/harm, and organisation turnover. & Definitive Guideline 2016 & Ch.29, Ch.30 \\
\midrule
Work-Related Deaths Protocol & Joint protocol between HSE, Police, CPS, CFOA, and LGA governing investigation of work-related deaths. & Protocol document; joint investigation procedures & Ch.29 \\
\bottomrule
\end{tabularx}
\end{table}
